\chapter{近平衡态中的输运过程}
\problem{3.1}{利用如图所示的旋转黏度计可以测定流体的黏度。已知由扭丝悬吊的两同轴圆筒的内筒的外半径和外筒的内半径分别为 $R$、$R+\delta$（$\delta\ll R$），长度为 $L$，待测流体装于两圆筒之间。使外圆筒以恒定角速度 $\omega$ 转动时，内圆筒会先随着转动，待转过一定角度后会“静止”下来，若该转角 $\theta$ 显示的扭力矩为 $M$，试证明待测的流体的黏度为
\[
\eta=\frac{M\delta}{2\pi R^2(R+\delta)L\omega}.
\]
\fig{3_1}}
\answerarea
\problem{3.2}{一细金属丝将一质量为 $m$、半径为 $R$ 的均质圆盘沿中心轴垂挂起来，并且盘能够轴转动，与盘面平行在其下置一大的水平板，盘与平板之间间距为 $d$，并充满黏度为 $\eta$ 的液体。初始时，盘以角速度 $\omega_0$ 旋转。假设盘与平板间的任一竖直线上的速度梯度都相等，试问在 $t$ 时刻盘的旋转角速度是多少？}
\answerarea
\clearpage
\problem{3.3}{一项观测表明，地球表面被晒热的地方形成了竖直向上的稳定气流，其速度为 $0.5\,\mathrm{m/s}$，并且在气流中有一球形尘埃以恒定速度 $0.1\,\mathrm{m/s}$ 向上运动。如果尘埃的密度为 $5.0\times10^3\,\mathrm{kg/m^3}$，空气的密度为 $1.3\,\mathrm{kg/m^3}$，空气的黏度为 $1.6\times10^{-5}\,\mathrm{Pa\cdot s}$，试确定尘埃的半径。}
\answerarea
\problem{3.4}{半径 $r=0.1\,\mathrm m$ 的铀球，在原子核裂变过程中以 $5.5\times10^3\,\mathrm{W/m^3}$ 的热产生率均匀地、恒定不变地散发出热量。已知铀的热导率为 $46\,\mathrm{W/(m\cdot K)}$，试问达到稳态时，铀球的中心与外表面间的温度差是多少？}
\answerarea
\clearpage
\problem{3.5}{厚度都可忽略的两个长圆筒共轴套在一起，两筒长度都为 $L$，内筒和外筒的半径分别为 $R_1,R_2$，并且两筒分别保持在恒定的温度 $T_1,T_2$（$T_1>T_2$）。设两筒间空气的热导率 $\kappa$ 对温度的依赖可忽略，试证明：每秒内由内筒通过空气传到外筒的热量为
\[
Q=\frac{2\pi\kappa L}{\ln(R_2/R_1)}(T_1-T_2).
\]}
\answerarea
\problem{3.6}{厚度为 $d$、面积为 $S$ 的平板，如果其两表面的温度分别为 $T_1,T_2$，形成该平板的物质的热导率与摄氏温标下的温度 $t$ 之间的关系为 $\kappa=\kappa_0(1+\alpha t)$，其中 $\kappa_0$ 和 $\alpha$ 为常量。
\begin{enumerate}
\item 试在忽略侧面热流的情况下，给出该平板的两表面间的热流；
\item 如果两表面的温度分别为 $T_1=0\,^\circ\mathrm C$、$T_2=100\,^\circ\mathrm C$，试确定该板正中间截面上的温度。
\end{enumerate}}
\answerarea
\clearpage
\problem{3.7}{现有尺寸相同、热导率相差一倍的两根金属棒，其侧面可以认为是绝热的。如果分别采用将它们串联和并联方式在温度分别恒定的高温处和低温处之间传热，两种连接方式的传热效率相差多少？}
\answerarea
\problem{3.8}{一人的体表温度大约为 $33\,^\circ\mathrm C$，正常衣物的热导率大约为 $3\,\mathrm{mW/(m\cdot K)}$，当其在气温为 $-13\,^\circ\mathrm C$ 的室外时，需要穿多厚的衣服才能感觉舒适？}
\answerarea
\clearpage
\problem{3.9}{在地面上方设置透明塑料薄膜，形成长 $50\,\mathrm m$、宽 $20\,\mathrm m$ 的温室大棚。设太阳表面的温度为 $5500\,^\circ\mathrm C$，太阳的半径为 $7.0\times10^8\,\mathrm m$，太阳与地球间的距离为 $1.5\times10^{11}\,\mathrm m$，并且太阳和地球都可以近似为绝对黑体；塑料薄膜的反射系数为 0.6，单从辐射角度讲，塑料薄膜也可以近似为黑体；大气对太阳能的吸收等的耗散率为 0.4。试问大棚内可能达到的最高温度为多少摄氏度，并将所得结果与没有薄膜时的结果进行比较。}
\answerarea
\problem{3.10}{设有一半径为 $R$ 的水滴悬浮在空气中，由于蒸发而体积逐渐减小，蒸发出的水蒸气扩散到周围空气中。设其邻近处水蒸气的密度为 $\rho$，远处水蒸气的密度为 $\rho_\infty$。水蒸气在空气中的扩散系数为 $D$，水的密度为 $\rho_W$，试证明：
\begin{enumerate}
\item 水滴的蒸发速率为 $W=4\pi D(\rho-\rho_\infty)R$；
\item 水滴全部蒸发完需要的时间为 $t=\dfrac{\rho_W R^2}{2D(\rho-\rho_\infty)}$。
\end{enumerate}}
\answerarea
\clearpage
\problem{3.11}{一长为 $2\,\mathrm m$、截面积为 $10^{-4}\,\mathrm{m^2}$ 的管子里贮有标准状态下的 $\mathrm{CO_2}$ 气体，一半 $\mathrm{CO_2}$ 分子中的 C 原子是放射性同位素 $^{14}\mathrm C$。在 $t=0$ 时，具有放射性的分子密集在管子左端，其分子数密度沿着管子均匀地减少，到右端减为 0。假设 $\mathrm{CO_2}$ 的黏度为 $1.4\times10^{-5}\,\mathrm{N\cdot s/m^2}$，试确定：
\begin{enumerate}
\item 开始时，放射性气体的密度梯度是多大？
\item 开始时，每秒有多少个放射性分子通过管子中点的横截面从左侧移往右侧？从右侧移往左侧的呢？
\item 开始时，每秒通过管子横截面扩散的放射性气体为多少克？
\end{enumerate}}
\answerarea
\problem{3.12}{在标准状态下，氦气的黏度为 $1.89\times10^{-5}\,\mathrm{Pa\cdot s}$，试确定此状态下氦原子的平均自由程和半径。}
\answerarea
\clearpage
\problem{3.13}{氧气在标准状态下的扩散系数为 $1.9\times10^{-5}\,\mathrm{m^2/s}$，试确定氧分子的平均自由程。}
\answerarea
\problem{3.14}{一定量的气体先经过等体过程，使其温度升高一倍；再经过等温过程，使其体积膨胀为原来的二倍。试问末态中气体分子的平均自由程 $\overline\lambda$、末态气体的黏度 $\eta$、热传导系数 $\kappa$、扩散系数 $D$ 各为原来的多少倍？}
\answerarea
\clearpage
\problem{3.15}{卫星飞行在极稀薄的大气中，这时大气分子的平均自由程比卫星的线度大得多。记大气分子的数密度为 $n$，大气分子质量为 $m$，卫星的飞行速率为 $v$，试确定气体作用在与卫星运动方向垂直的单位面积器壁上的阻力。}
\answerarea
\problem{3.16}{在稀薄气体中平行放置两片薄板 $A,B$，这两片薄板分别以速率 $v_A,v_B$ 沿平行板面方向运动。记气体的密度为 $\rho$，气体分子的平均速率为 $\overline v$，试证明：作用在板上单位面积的黏滞力为
\[
F=\frac16\rho\overline v(v_A-v_B).
\]}
\answerarea
\clearpage
\problem{3.17}{在稀薄气体中平行放置温度分别为 $T_A,T_B$ 的薄板 $A,B$。记气体的密度为 $\rho$，气体分子的平均速率为 $\overline v$，气体的定容比热容为 $C_V$，试证明：因两板之间有温度差，在单位时间内通过单位面积传递的热量为
\[
Q=\frac16\rho\overline v C_V(T_A-T_B).
\]}
\answerarea
\problem{3.18}{在 $18\,^\circ\mathrm C$ 的温度下，观察半径为 $0.4\,\mu\mathrm m$ 的粒子在黏度为 $2.78\times10^{-3}\,\mathrm{Pa\cdot s}$ 的液体中的布朗运动。测得粒子在时间间隔 $10\,\mathrm s$ 内的位移方差为 $\overline{x^2}=3.3\times10^{-12}\,\mathrm{m^2}$。试确定玻尔兹曼常量 $k_{\rm B}$。}
\answerarea
\clearpage
